% page 494 du fac-simile (image p-493 du scan)
% bloc 160.0 x 237.7 mm ; marge gauche 8.0 mm
\pgc{-12.700mm}{0.000mm}{
\l{\cel{3}486}
\l{}
\l{\sou{referar}. (\sou{trans.,}\cel{1}ad). I. Relatigar (ulo) a to quo explikas}
\l{\cel{3}lu, konfirmas lu. - II.\cel{1}\sou{Referar su}\cel{1}ad :\cel{2}rekursar a la au-}
\l{\cel{3}toritato di ulu. - DEFIRS.}
\l{}
\l{\sou{referendario}. Oficiero di kancelereyo qua havas kom taski}
\l{\cel{3}gardar la rejo-siglilo e raportar pri la supliko-letri.}
\l{\cel{3}- DEF.}
\l{}
\l{\sou{reflektar}. I. (\sou{trans.)}\cel{2}(\sou{aludante surfaco}). Retrosendar (en}
\l{\cel{3}ula kazi : plu feble) to quo frapas lu. - II. (\sou{netrans}.,}
\l{\cel{3}\sou{pri}). Retrovenar a sua pensajo por studiar lu plu profunde.}
\l{\cel{3}- DEFIRS.}
\l{}
\l{\sou{reflektoro.}\cel{2}(\sou{tekn., optik}o).Aparato qua konsisas ye surfaco}
\l{\cel{3}de metalo polisita od ye spegulo, destinita reflektor la}
\l{\cel{3}radii konsiderata segun la direciono selektita ; li uzesas}
\l{\cel{3}precipue por lumo, ma anke por soni, kaloro. - DEFS.}
\l{}
\l{\sou{refluxar}. (\sou{netrans.}) (\sou{pri mareo}). Decensar, pos acensir per}
\l{\cel{3}fluxo. - EFIS.}
\l{}
\l{\sou{reflexa}. (\sou{fiziol.}) Produktita sen-vole par reakto organisma-}
\l{\cel{3}la. - DEFIRS.}
\l{}
\l{\sou{reformar}. (\sou{trans.}) I. (\sou{eklezio kristan}.)Retroduktar a la for-}
\l{\cel{3}mo originala. - II. Donar formo plu bona, qua igas plu}
\l{\cel{3}apta por la skopo. - DEFIRS.}
\l{}
\l{\sou{refraktar}. (\sou{trans.)}\cel{2}(\sou{fiziko}).Igar (lumo-radio) deviacar pro}
\l{\cel{3}(per) lua paso (pasigo) en medio di qua la denseso-gradi}
\l{\cel{3}rispektiva diferas. - DEFIRS.}
\l{}
\l{\sou{refraktara}.\cel{2}I. (\sou{tekn.)}\cel{2}Qua rezistas ula influi. - II.(\sou{milit}.)}
\l{\cel{3}Qua eskapas la lego pri armeanigado. - DEFIRS.}
\l{}
\l{\sou{refraktoro}. (\sou{tekn.)}\cel{2}Instrumento qua donas la indici di re-}
\l{\cel{3}frakto. - DEFIRS.}
\l{}
\l{refreno. I. Ri-aparo di verso en omna kupleto di kansono, di}
\l{\cel{3}balado, di rondelo. - II. to quon onu repetas sencese.-}
\l{}
\l{}
\l{\sou{refujar}. (\sou{netrans., en, che)}.\cel{2}Irar rapide a loko ube onu po-}
\l{\cel{3}vas vivar sekure. - EFIS.}
\l{}
\l{\sou{refutar}. (\sou{trans.}) Destruktar la valideso (di opiniono) per}
\l{\cel{3}demonstrar ke lu esas ne-exakta. - DEFIS.}
\l{}
\l{\sou{refuzar}. (\sou{trans.)}\cel{2}I. Ne-grantar. - II. Ne-aceptar. - III.}
\l{\cel{3}Ne-admisar. - EFIS.}
\l{}
\l{\sou{regala}r (\sou{trans., per)}. I. Restorar (ulu) per ulo bona drinkale}
\l{\cel{3}o manjale. - II. Ofrar (a gasti) partio di plezuro. - DEFS.}
}
